% GENERATED FILE. Edit canonical lessons, then run npm run generate:books.
% canonical-source-hash: fnv1a64:a40261d39f5b786c
% canonical-lessons: GU-C454-ucat, GU-C454-phikar, GU-C454-becen, GU-C454-vyakul, GU-C454-salvalvu, GU-R454-review-447-454

\chapter{A Restless Night}
\label{ch:gu-a-restless-night}
\hlchaptermodality{454}

\begin{chapteropening}
\textbf{By the end of this chapter:} \emph{I can say what keeps me awake with worry.}

Complete the last lesson of chapter 454: I can say what keeps me awake with worry.
\end{chapteropening}

\section[\texorpdfstring{\gu{ઉચાટ} (ucāṭ)}{ઉચાટ (ucāṭ)}]{\gu{ઉચાટ} (ucāṭ) — anxious unease}

\label{lesson:GU-C454-ucat}



\begin{quote}
Before the new one: say the Gujarati for torment, then the Gujarati for apprehension.
\end{quote}

\subsection*{You'll want to know: \gu{ઉચાટ}}
\textbf{\gu{ઉચાટ}} — \emph{ucāṭ} — \textquotedblleft{}anxious unease\textquotedblright{}.

\begin{quote}
\gu{દીકરો} \gu{મોડે} \gu{સુધી} \gu{ઘરે} \gu{ન} \gu{આવ્યો}, \gu{તેથી} \gu{માતાને} \gu{ઉચાટ} \gu{થયો}.
\end{quote}

\begin{quote}
\emph{dīkro moḍe sudhī ghare na āvyo, tethī mātāne ucāṭ thayo.}
\end{quote}

\begin{quote}
The son had not come home till late, so his mother grew anxious.
\end{quote}

A masculine noun.

\begin{cousinweb}
\textbf{\gu{ઉચાટ}} (\emph{ucāṭ}) is inherited: CDIAL lists the Gujarati word, \textquotedblleft{}impatience, anxiety\textquotedblright{}, under Sanskrit \emph{uccaṭati}, \textquotedblleft{}goes away\textquotedblright{} (CDIAL 1635): a mind that will not stay put.
\end{cousinweb}

\begin{grammarlens}[title={What you've built}]
One more word for describing what you lived and felt.
\end{grammarlens}

\subsection*{Your turn}
\begin{itemize}
\raggedright
  \item \emph{Say it:} \emph{ucāṭ}
  \item \emph{Say it:} \emph{ucāṭ}, once more
  \item \emph{Say it:} \emph{dīkro moḍe sudhī ghare na āvyo, tethī mātāne ucāṭ thayo.} — The son had not come home till late, so his mother grew anxious.
  \item \emph{Recall:} say the Gujarati for torment, then the Gujarati for apprehension, then say \emph{ucāṭ} again
\end{itemize}

\begin{tcolorbox}[breakable,colback=teal!4,colframe=teal!35!black,title={Before you move on}]
What does \gu{ઉચાટ} mean? (\textquotedblleft{}Anxious unease\textquotedblright{}.) Say it once more.
\end{tcolorbox}
\section[\texorpdfstring{\gu{ફિકર} (phikar)}{ફિકર (phikar)}]{\gu{ફિકર} (phikar) — concern; worry about someone}

\label{lesson:GU-C454-phikar}



\begin{quote}
Before the new one: say the Gujarati for apprehension, then the Gujarati for anxious unease.
\end{quote}

\subsection*{You'll want to know: \gu{ફિકર}}
\textbf{\gu{ફિકર}} — \emph{phikar} — \textquotedblleft{}concern; worry about someone\textquotedblright{}.

\begin{quote}
\gu{તમે} \gu{મારી} \gu{ફિકર} \gu{ન} \gu{કરો}.
\end{quote}

\begin{quote}
\emph{tame mārī phikar na karo.}
\end{quote}

\begin{quote}
Don't worry about me.
\end{quote}

A feminine noun.

\begin{cousinweb}
\textbf{\gu{ફિકર}} (\emph{phikar}) is usually traced to Arabic \emph{fikr}, \textquotedblleft{}thought, reflection\textquotedblright{}, which reached Gujarati through Persian with the sense \textquotedblleft{}worry\textquotedblright{}.
\end{cousinweb}

\begin{grammarlens}[title={What you've built}]
One more word for describing what you lived and felt.
\end{grammarlens}

\subsection*{Your turn}
\begin{itemize}
\raggedright
  \item \emph{Say it:} \emph{phikar}
  \item \emph{Say it:} \emph{phikar}, once more
  \item \emph{Say it:} \emph{tame mārī phikar na karo.} — Don't worry about me.
  \item \emph{Recall:} say the Gujarati for apprehension, then the Gujarati for anxious unease, then say \emph{phikar} again
\end{itemize}

\begin{tcolorbox}[breakable,colback=teal!4,colframe=teal!35!black,title={Before you move on}]
What does \gu{ફિકર} mean? (\textquotedblleft{}Concern; worry about someone\textquotedblright{}.) Say it once more.
\end{tcolorbox}
\section[\texorpdfstring{\gu{બેચેન} (becen)}{બેચેન (becen)}]{\gu{બેચેન} (becen) — restless; uneasy}

\label{lesson:GU-C454-becen}



\begin{quote}
Before the new one: say the Gujarati for anxious unease, then the Gujarati for concern.
\end{quote}

\subsection*{You'll want to know: \gu{બેચેન}}
\textbf{\gu{બેચેન}} — \emph{becen} — \textquotedblleft{}restless; uneasy\textquotedblright{}.

\begin{quote}
\gu{પરિણામ} \gu{પહેલાં} \gu{મીરા} \gu{બેચેન} \gu{હતી}.
\end{quote}

\begin{quote}
\emph{pariṇām pahelā̃ Mīrā becen hatī.}
\end{quote}

\begin{quote}
Mira was uneasy before the result.
\end{quote}

\begin{cousinweb}
\textbf{\gu{બેચેન}} (\emph{becen}) is usually traced to Persian \emph{bē-čain}, \textquotedblleft{}without rest\textquotedblright{}, from \emph{bē}, \textquotedblleft{}without\textquotedblright{}, and \emph{čain}, \textquotedblleft{}rest, ease\textquotedblright{}.
\end{cousinweb}

\begin{grammarlens}[title={What you've built}]
One more word for describing what you lived and felt.
\end{grammarlens}

\subsection*{Your turn}
\begin{itemize}
\raggedright
  \item \emph{Say it:} \emph{becen}
  \item \emph{Say it:} \emph{becen}, once more
  \item \emph{Say it:} \emph{pariṇām pahelā̃ Mīrā becen hatī.} — Mira was uneasy before the result.
  \item \emph{Recall:} say the Gujarati for anxious unease, then the Gujarati for concern, then say \emph{becen} again
\end{itemize}

\begin{tcolorbox}[breakable,colback=teal!4,colframe=teal!35!black,title={Before you move on}]
What does \gu{બેચેન} mean? (\textquotedblleft{}Restless; uneasy\textquotedblright{}.) Say it once more.
\end{tcolorbox}
\section[\texorpdfstring{\gu{વ્યાકુળ} (vyākuḷ)}{વ્યાકુળ (vyākuḷ)}]{\gu{વ્યાકુળ} (vyākuḷ) — distraught}

\label{lesson:GU-C454-vyakul}



\begin{quote}
Before the new one: say the Gujarati for concern, then the Gujarati for restless.
\end{quote}

\subsection*{You'll want to know: \gu{વ્યાકુળ}}
\textbf{\gu{વ્યાકુળ}} — \emph{vyākuḷ} — \textquotedblleft{}distraught\textquotedblright{}.

\begin{quote}
\gu{બાળક} \gu{ખોવાઈ} \gu{ગયું} \gu{ત્યારે} \gu{માતા} \gu{વ્યાકુળ} \gu{થઈ} \gu{ગઈ}.
\end{quote}

\begin{quote}
\emph{bāḷak khovāī gayũ tyāre mātā vyākuḷ thaī gaī.}
\end{quote}

\begin{quote}
When the child got lost, the mother was distraught.
\end{quote}

\begin{cousinweb}
\textbf{\gu{વ્યાકુળ}} (\emph{vyākuḷ}) is a tatsama, Sanskrit \emph{vyākula}, \textquotedblleft{}entirely filled with, confused, bewildered\textquotedblright{}, as Monier-Williams records it.
\end{cousinweb}

\begin{grammarlens}[title={What you've built}]
One more word for describing what you lived and felt.
\end{grammarlens}

\subsection*{Your turn}
\begin{itemize}
\raggedright
  \item \emph{Say it:} \emph{vyākuḷ}
  \item \emph{Say it:} \emph{vyākuḷ}, once more
  \item \emph{Say it:} \emph{bāḷak khovāī gayũ tyāre mātā vyākuḷ thaī gaī.} — When the child got lost, the mother was distraught.
  \item \emph{Recall:} say the Gujarati for concern, then the Gujarati for restless, then say \emph{vyākuḷ} again
\end{itemize}

\begin{tcolorbox}[breakable,colback=teal!4,colframe=teal!35!black,title={Before you move on}]
What does \gu{વ્યાકુળ} mean? (\textquotedblleft{}Distraught\textquotedblright{}.) Say it once more.
\end{tcolorbox}
\section[\texorpdfstring{\gu{સળવળવું} (saḷvaḷvũ)}{સળવળવું (saḷvaḷvũ)}]{\gu{સળવળવું} (saḷvaḷvũ) — to stir; to fidget}

\label{lesson:GU-C454-salvalvu}



\begin{quote}
Before the new one: say the Gujarati for restless, then the Gujarati for distraught.
\end{quote}

\subsection*{You'll want to know: \gu{સળવળવું}}
\textbf{\gu{સળવળવું}} — \emph{saḷvaḷvũ} — \textquotedblleft{}to stir; to fidget\textquotedblright{}.

\begin{quote}
\gu{ચિંતાને} \gu{લીધે} \gu{અરુણ} \gu{આખી} \gu{રાત} \gu{પથારીમાં} \gu{સળવળતો} \gu{હતો}.
\end{quote}

\begin{quote}
\emph{cintāne līdhe Aruṇ ākhī rāt pathārīmā̃ saḷvaḷto hato.}
\end{quote}

\begin{quote}
Arun was tossing and fidgeting in bed all night with worry.
\end{quote}

Past (he …): \gu{તે} \gu{પથારીમાં} \gu{સળવળ્યો}, \emph{te pathārīmā̃ saḷvaḷyo}.

\begin{cousinweb}
\textbf{\gu{સળવળવું}} (\emph{saḷvaḷvũ}) is inherited from Sanskrit \emph{salati}, \textquotedblleft{}moves\textquotedblright{}; CDIAL lists the Gujarati verb beside \emph{saḷvaḷ}, \textquotedblleft{}writhing about\textquotedblright{} (CDIAL 13284).
\end{cousinweb}

\begin{grammarlens}[title={What you've built}]
One more word for describing what you lived and felt.
\end{grammarlens}

\subsection*{Your turn}
\begin{itemize}
\raggedright
  \item \emph{Say it:} \emph{saḷvaḷvũ}
  \item \emph{Say it:} \emph{saḷvaḷvũ}, once more
  \item \emph{Say it:} \emph{cintāne līdhe Aruṇ ākhī rāt pathārīmā̃ saḷvaḷto hato.} — Arun was tossing and fidgeting in bed all night with worry.
  \item \emph{Recall:} say the Gujarati for restless, then the Gujarati for distraught, then say \emph{saḷvaḷvũ} again
\end{itemize}

\begin{tcolorbox}[breakable,colback=teal!4,colframe=teal!35!black,title={Before you move on}]
What does \gu{સળવળવું} mean? (\textquotedblleft{}To stir; to fidget\textquotedblright{}.) Say it once more.
\end{tcolorbox}
\section[(dialogue)]{Looking back: from A Morning Full of Joy to A Restless Night}

\label{lesson:GU-R454-review-447-454}



\begin{quote}
Say the Gujarati for delight and to rejoice.
\end{quote}

\subsection*{Your turn}
Say these aloud:

\begin{itemize}
\raggedright
  \item delight, to rejoice, warmth of feeling, to experience, to delight in
  \item laughter, to make someone laugh, a smile, a joke, to tease
  \item eager, keen, longing, to be won over, pleased
  \item fondness, affection, a mother's love, to captivate, pampering
  \item to cling, to kiss, to hold on tight, love, to clasp
  \item to get cross, wrath, temper, to seethe, strife
  \item to be scared, awe, to shudder, torment, apprehension
  \item anxious unease, concern, restless, distraught, to fidget
\end{itemize}

\begin{tcolorbox}[breakable,colback=teal!4,colframe=teal!35!black,title={Before you move on}]
Delight? (\textbf{\gu{હરખ}}, \emph{harakh}.) To rejoice? (\textbf{\gu{હરખાવું}}, \emph{harkhāvũ}.) Warmth of feeling? (\textbf{\gu{ઉમળકો}}, \emph{umaḷko}.) To experience? (\textbf{\gu{અનુભવવું}}, \emph{anubhavvũ}.) To delight in? (\textbf{\gu{રાચવું}}, \emph{rācvũ}.) Laughter? (\textbf{\gu{હાસ્ય}}, \emph{hāsya}.) To make someone laugh? (\textbf{\gu{હસાવવું}}, \emph{hasāvvũ}.) A smile? (\textbf{\gu{સ્મિત}}, \emph{smit}.) A joke? (\textbf{\gu{મજાક}}, \emph{majāk}.) To tease? (\textbf{\gu{ચીડવવું}}, \emph{cīḍavvũ}.) Eager? (\textbf{\gu{આતુર}}, \emph{ātur}.) Keen? (\textbf{\gu{ઉત્સુક}}, \emph{utsuk}.) Longing? (\textbf{\gu{ઉત્કંઠા}}, \emph{utkanṭhā}.) To be won over? (\textbf{\gu{રીઝવું}}, \emph{rījhvũ}.) Pleased? (\textbf{\gu{પ્રસન્ન}}, \emph{prasann}.) Fondness? (\textbf{\gu{વહાલ}}, \emph{vahāl}.) Affection? (\textbf{\gu{સ્નેહ}}, \emph{sneh}.) A mother's love? (\textbf{\gu{મમતા}}, \emph{mamtā}.) To captivate? (\textbf{\gu{મોહવું}}, \emph{mohvũ}.) Pampering? (\textbf{\gu{લાડ}}, \emph{lāḍ}.) To cling? (\textbf{\gu{વળગવું}}, \emph{vaḷagvũ}.) To kiss? (\textbf{\gu{ચૂમવું}}, \emph{cūmvũ}.) To hold on tight? (\textbf{\gu{બાઝવું}}, \emph{bājhvũ}.) Love? (\textbf{\gu{પ્યાર}}, \emph{pyār}.) To clasp? (\textbf{\gu{ચાંપવું}}, \emph{cā̃pvũ}.) To get cross? (\textbf{\gu{ખીજાવું}}, \emph{khījāvũ}.) Wrath? (\textbf{\gu{ક્રોધ}}, \emph{krodh}.) Temper? (\textbf{\gu{મિજાજ}}, \emph{mijāj}.) To seethe? (\textbf{\gu{ઉકળવું}}, \emph{ukaḷvũ}.) Strife? (\textbf{\gu{કલહ}}, \emph{kalah}.) To be scared? (\textbf{\gu{બીવું}}, \emph{bīvũ}.) Awe? (\textbf{\gu{ધાક}}, \emph{dhāk}.) To shudder? (\textbf{\gu{થથરવું}}, \emph{thatharvũ}.) Torment? (\textbf{\gu{ત્રાસ}}, \emph{trās}.) Apprehension? (\textbf{\gu{ભીતિ}}, \emph{bhīti}.) Anxious unease? (\textbf{\gu{ઉચાટ}}, \emph{ucāṭ}.) Concern? (\textbf{\gu{ફિકર}}, \emph{phikar}.) Restless? (\textbf{\gu{બેચેન}}, \emph{becen}.) Distraught? (\textbf{\gu{વ્યાકુળ}}, \emph{vyākuḷ}.) To fidget? (\textbf{\gu{સળવળવું}}, \emph{saḷvaḷvũ}.)
\end{tcolorbox}
